
\section{Background}

\subsection{Dense Softmax Attention}

For $Q, K, V \in \mathbb{R}^{L \times d_h}$:

\begin{equation}
\mathcal{A}_\text{dense}(Q, K, V) = \text{softmax}\!\left(\frac{Q K^\top}{\sqrt{d_h}}\right) V \in \mathbb{R}^{L \times d_h}
\end{equation}

\textbf{Complexity.} $\mathcal{O}(L^2 d_h)$ time, $\mathcal{O}(L^2)$ memory.

\textbf{Causality.} For decoder models, mask: $S_{ij} = -\infty$ for $j > i$.

\subsection{Linear Attention}

Replace softmax with a kernel $\phi: \mathbb{R}^{d_h} \to \mathbb{R}^m$:

\begin{equation}
\mathcal{A}_\text{lin}(Q, K, V) = \frac{\phi(Q) \left(\phi(K)^\top V\right)}{\phi(Q) \left(\phi(K)^\top \mathbf{1}\right)}
\end{equation}

where $\mathbf{1} \in \mathbb{R}^L$ is a vector of ones.

\textbf{Complexity.} $\mathcal{O}(L m d_h)$ time, $\mathcal{O}(m d_h)$ memory.

\textbf{Non-causal.} Direct computation via $\phi(K)^\top V \in \mathbb{R}^{m \times d_h}$.

\textbf{Causal.} Use prefix sums:
\begin{equation}
\text{Num}_i = \phi(Q_i) \cdot \sum_{j \le i} \phi(K_j) V_j^\top, \quad
\text{Den}_i = \phi(Q_i) \cdot \sum_{j \le i} \phi(K_j)
\end{equation}

\subsection{Rectified Kernel}

We use $\phi(x) = \mathrm{elu}(x) + 1$ where:

\begin{equation}
\mathrm{elu}(x) = \begin{cases} x & x > 0 \\ \alpha (e^x - 1) & x \leq 0 \end{cases}
\end{equation}

with $\alpha = 1$. This kernel:

\begin{itemize}
\item Is \textbf{non-negative}: $\phi(x) > 0$ for all $x$.
\item Is \textbf{continuous} and \textbf{smooth}.
\item Is \textbf{bounded} in expectation for bounded inputs.

\end{itemize}
\textbf{Why non-negativity matters.} In (3.2), the denominator 
$\phi(Q) (\phi(K)^\top \mathbf{1})$ is a sum of products of non-negative 
terms, hence positive. This avoids division by zero or negative denominators.

\subsection{Prefix Sum for Causal Attention}

For causal linear attention, we compute:

\begin{equation}
\text{cum}_i = \sum_{j \le i} \phi(K_j) V_j^\top \in \mathbb{R}^{m \times d_h}
\end{equation}

This is a \textbf{cumulative sum} over the sequence dimension, which we implement 
via \texttt{torch.cumsum}. The complexity is $\mathcal{O}(L m d_h)$ memory and time.

\textbf{Numerical stability.} Cumsum in FP32 avoids overflow for $L \le 8192$ 
with $m \le 128$. For longer sequences, chunked cumsum or a recurrent 
formulation is needed.

\subsection{Gating}

We use a per-head sigmoid gate:

\begin{equation}
g_{h,i} = \sigma\!\left(w_h^\top x_i + b_h\right) \in (0, 1)
\end{equation}

The output is a convex combination:

\begin{equation}
O_i = g_{h,i} \cdot O_i^\text{dense} + (1 - g_{h,i}) \cdot O_i^\text{lin}
\end{equation}

\textbf{Interpretation.} $g \approx 1$ $\to$ dense attention; $g \approx 0$ $\to$ linear 
attention; intermediate $\to$ smooth mix.

\subsection{Grouped Query Attention (GQA)}

GQA (Ainslie et al., 2023) reduces $H_{kv} < H$ and repeats KV heads:

\begin{equation}
K, V \in \mathbb{R}^{B \times H_{kv} \times L \times d_h}, \quad
\tilde{K}, \tilde{V} = \text{repeat\_kv}(K, V, H / H_{kv})
\end{equation}

HALO-ReGLA supports GQA via \texttt{repeat\_kv} before attention computation.

\subsection{Rotary Position Embeddings (RoPE)}

RoPE (Su et al., 2021) applies position-dependent rotation:

\begin{equation}
\tilde{Q}_i = Q_i \cos(\theta_i) + \mathrm{rot}(Q_i) \sin(\theta_i)
\end{equation}

HALO-ReGLA applies RoPE to Q, K before both branches.

\subsection{LoRA Recovery}

LoRA (Hu et al., 2021) adds trainable adapters:

\begin{equation}
W_\text{eff} = W_\text{frozen} + \frac{\alpha}{r} B A
\end{equation}

with $B = 0$ at initialization. HALO-ReGLA uses LoRA to adapt the frozen 
projections.

\subsection{Notation Summary}

\begin{lstlisting}

---
\end{lstlisting}